\usepackage{xeCJK}
\usepackage{geometry}
\geometry{left=2cm,right=2cm,top=2.5cm,bottom=2.5cm}
\setlength{\headheight}{12.7pt}

\usepackage{titlesec}
\titleformat{\section}
  {\normalfont\large\bfseries}
  {\thesection}
  {1em}
  {}
\titlespacing*{\section}
  {0pt}
  {3.5ex plus 1ex minus .2ex}
  {2.3ex plus .2ex}

\usepackage{amsmath}
\usepackage{amssymb}
\usepackage{amsthm}
\usepackage{enumitem}
\setlist[enumerate]{itemsep=0pt,topsep=0pt,parsep=0pt,leftmargin=0pt,itemindent=2.5em}
\setlist[itemize]{itemsep=0pt,topsep=0pt,parsep=0pt,leftmargin=0pt,itemindent=2.5em}
\usepackage{fancyhdr}
\pagestyle{fancy}
\fancyhf{}
\usepackage{graphicx}
\usepackage{hyperref}
\usepackage{indentfirst}
\usepackage{lmodern}


\begin{document}
\setlength{\abovedisplayskip}{1pt}
\setlength{\belowdisplayskip}{1pt}

\section{语言的限制与扩张}

\hypertarget{sec:定义1.1}{}
\noindent\textbf{定义1.1}

给定一阶语言$\mathcal{L}_1,\mathcal{L}_2$, 设$\mathcal{L}_1\subset\mathcal{L}_2$,
再给定$\mathcal{L}_1$-结构$A_1$和$\mathcal{L}_2$-结构$A_2$.

如果$A_1,A_2$有相同的论域$A$, 并且$\mathcal{L}_1$内所有符号在$A_1,A_2$中有相同的解释,
就称$A_1$是$A_2$的\textbf{解释限制}, 或称$A_2$是$A_1$的\textbf{解释扩张}.

\vspace{10pt}

在此设定下, $\mathcal{L}_1$的项和公式在$\mathcal{L}_2$中有同样的语义,
具体有下面关于语义的定理.

\vspace{10pt}

\hypertarget{sec:定理1.1}{}
\noindent\textbf{定理1.1}

任给$\sigma_1:\mathcal{V}_1\to A,\sigma_2:\mathcal{V}_2\to A$
和$\mathcal{L}_1$的项$t$, 如果$\sigma_1=\sigma_2\!\restriction_{\mathcal{V}_1}$, 那么
$t^{(A_1,\sigma_1)}=t^{(A_2,\sigma_2)}$.

\noindent{\kaishu 证明.}

\begin{enumerate}
    \item 对任意的常元$c$, $c^{(A_1,\sigma_1)}=c^{A_1}=c^{A_2}=c^{(A_2,\sigma_2)}$,
    \item 对任意的变元$x$,
    $x^{(A_1,\sigma_1)}=x^{\sigma_1}=x^{\sigma_2}=x^{(A_2,\sigma_2)}$,
    \item 对任意的$n$元函数符号$f$, 假设项$t_1,\cdots,t_n$满足命题, 那么
    \begin{equation*}
        (ft_1\cdots t_n)^{(A_1,\sigma_1)}
        =f^{A_1}(t_1^{(A_1,\sigma_1)},\cdots,t_n^{(A_1,\sigma_1)})
        =f^{A_2}(t_1^{(A_2,\sigma_2)},\cdots,t_n^{(A_2,\sigma_2)})
        =(ft_1\cdots t_n)^{(A_2,\sigma_2)}.
    \tag*{\qedsymbol}\end{equation*}
\end{enumerate}

\vspace{10pt}

\hypertarget{sec:定理1.2}{}
\noindent\textbf{定理1.2}

任给$\sigma_1:\mathcal{V}_1\to A,\sigma_2:\mathcal{V}_2\to A$
和$\mathcal{L}_1$的公式$\varphi$, 如果$\sigma_1=\sigma_2\!\restriction_{\mathcal{V}_1}$, 那么
$(A_1,\sigma_1)\models^1\varphi\leftrightarrow(A_2,\sigma_2)\models^2\varphi$.

\noindent{\kaishu 证明.}

\begin{enumerate}
    \item 对任意的项$t_1,t_2$,
    \[
        (A_1,\sigma_1)\models^1t_1\equiv t_2\iff
        t_1^{(A_1,\sigma_1)}=t_2^{(A_1,\sigma_1)}\iff
        t_1^{(A_2,\sigma_2)}=t_2^{(A_2,\sigma_2)}\iff
        (A_2,\sigma_2)\models^2t_1\equiv t_2,
    \]
    \item 对任意的$n$元关系符号$r$和项$t_1,\cdots,t_n$,
    \[\begin{aligned}
        &(A_1,\sigma_1)\models^1rt_1\cdots t_n\iff
        (t_1^{(A_1,\sigma_1)},\cdots,t_n^{(A_1,\sigma_1)})\in r^{A_1}\\[-3pt]
        \iff&(t_1^{(A_2,\sigma_2)},\cdots,t_n^{(A_2,\sigma_2)})\in r^{A_2}\iff
        (A_2,\sigma_2)\models^2rt_1\cdots t_n,
    \end{aligned}\]
    \item 假设$\varphi_1$满足命题, 显然$\varphi=\neg\varphi_1$满足命题,
    \item 假设$\varphi_1,\varphi_2$满足命题,
    显然$\varphi=(\varphi_1\to\varphi_2)$满足命题,
    \item 对任意的变元$x$, 假设$\varphi_1$满足命题,
    那么对$\varphi=\exists x\varphi_1$有:
    \begin{equation*}
        (A_1,\sigma_1)\models^1\varphi\iff
        \exists a\in A:\left(A_1,\sigma_1\frac{a}{x}\right)\models^1\varphi_1\iff
        \exists a\in A:\left(A_2,\sigma_2\frac{a}{x}\right)\models^2\varphi_1\iff
        (A_2,\sigma_2)\models^2\varphi.
    \tag*{\qedsymbol}\end{equation*}
\end{enumerate}

\vspace{20pt}





\section{结构扩充语言}

\hypertarget{sec:定义2.1}{}
\noindent\textbf{定义2.1 (扩充语言)}

给定一阶语言$\mathcal{L}$和$\mathcal{L}$-结构$A$, 定义
\[
    \mathcal{L}_A\overset{\mathrm{def}}{=}(\mathcal{V},\mathcal{C}\cup A,
    (n_f)_{\mathcal{F}},(n_r)_{\mathcal{R}}),
\]
称之为结构$A$对语言$\mathcal{L}$的\textbf{扩充语言}.

\noindent{\kaishu 备注.
这里默认$A$与$\mathcal{L}$各集合均不相交, 不再展开讨论.}

\vspace{10pt}

进一步, $A$中的常元可以把自身当作它在$A$中的解释,
由此可以定义$A$在$\mathcal{L}_A$上的解释扩张$A^*$.

% 这样的定义未免太不严谨了.
% 但是严谨的写法非常的繁杂、无聊, 并且偏离了逻辑本身的意义.
% 希望未来使用这些概念时, 不会因为不严谨性而造成困扰.

\end{document}